% Portafolio final | FIQ-BUAP
\documentclass[11pt,letterpaper]{article}

\usepackage{iftex}
\ifPDFTeX
  \usepackage[T1]{fontenc}
  \usepackage[utf8]{inputenc}
  \usepackage{lmodern}
\else
  \usepackage{fontspec}
  \IfFontExistsTF{Libertinus Serif}
    {\setmainfont{Libertinus Serif}}{\setmainfont{Latin Modern Roman}}
  \IfFontExistsTF{Libertinus Sans}
    {\setsansfont{Libertinus Sans}}{\setsansfont{Latin Modern Sans}}
\fi
\usepackage{babel}
\babelprovide[main,import]{spanish}
\usepackage[margin=0.78in,top=0.83in,bottom=0.82in,
            headheight=17pt,headsep=11pt,footskip=27pt]{geometry}
\usepackage{graphicx}
\usepackage[table]{xcolor}
\usepackage{array,tabularx,longtable,booktabs,caption}
\usepackage{enumitem}
\usepackage{titlesec,fancyhdr,lastpage,pdflscape}
\usepackage{tcolorbox}
\usepackage{etoolbox,xparse}
\usepackage{hyperref}
\usepackage{bookmark}

\definecolor{Grafito}{HTML}{293333}
\definecolor{Terracota}{HTML}{C8715B}
\definecolor{Enlace}{HTML}{8E3C2D}
\definecolor{Arena}{HTML}{FAF4EC}
\definecolor{Tinta}{HTML}{232B2B}
\definecolor{Gris}{HTML}{526060}
\definecolor{Papel}{HTML}{F5F7F6}
\definecolor{Separador}{HTML}{D7DEDB}

\hypersetup{unicode,colorlinks=true,linkcolor=Enlace,urlcolor=Enlace,
  citecolor=Enlace,pdfstartview=FitH,bookmarksopen=true,bookmarksopenlevel=1}
\setlength{\parindent}{0pt}
\setlength{\parskip}{0.36em}
\setlength{\emergencystretch}{2em}
\setlength{\LTpre}{6pt}
\setlength{\LTpost}{7pt}
\setlist[itemize]{leftmargin=1.35em,itemsep=2pt,topsep=3pt}
\renewcommand{\arraystretch}{1.21}
\newcolumntype{Y}{>{\raggedright\arraybackslash}X}
\newcolumntype{L}[1]{>{\raggedright\arraybackslash}p{#1}}
\setcounter{tocdepth}{1}
\raggedbottom

% A. DATOS DEL EQUIPO
\newcommand{\Institucion}{Benemérita Universidad Autónoma de Puebla}
\newcommand{\Facultad}{Facultad de Ingeniería Química}
\newcommand{\Asignatura}{Diseño de Procesos y Productos Químicos}
\newcommand{\Docente}{Dr. Jesús Andrés Arzola Flores}
\newcommand{\Proyecto}{Planta de Producción de e-Metanol Grado Químico AA (30,000 t/año)}
\newcommand{\Firma}{Solvanta Engineering Solutions S.A. de C.V.}
\newcommand{\SiglasFirma}{Solvanta}
\newcommand{\FechaConformacion}{17/08/2026}
\newcommand{\Grupo}{1}
\newcommand{\Periodo}{Otoño 2026}
\newcommand{\FechaFinal}{23/11/2026}
\newcommand{\Revision}{Rev.\ 1.0}
\newcommand{\LogoEmpresa}{logo-empresa.png} % También admite PNG o JPG.

% Cada línea agrega un integrante a la tabla de portada.
\newcommand{\IntegrantesFirma}{%
  {Karol Paola Camacho López} & Directora de Proyecto / Optimización (Python) & Dirección ejecutiva, control documental, optimización SOEC y notebooks \\
  {Gildardo Ocaña Barrales} & Gerente de Procesos / DTI y P\&ID & Diagramación en Plant 3D (BFD/DFP/DTI), dimensionamiento y arreglo general \\
  {Paulina Mota Macip} & Ing. de Simulación (DWSIM) / Control & Modelado riguroso DWSIM, balances, tablas de proceso (TC/TE) y filosofía de control \\
  {Grecia Figueroa Rincón} & Especialista en HSE / Evaluación Económica & Evaluación ambiental (AC-AMB-01), HAZOP/SST, modelo BioSTEAM y evaluación financiera \\
}

\pagestyle{fancy}
\fancyhf{}
\fancyhead[L]{\small\sffamily\bfseries\color{Grafito} FIQ · BUAP}
\fancyhead[R]{\small\sffamily\color{Gris}\Asignatura}
\fancyfoot[L]{\small\sffamily\color{Grafito}\Docente}
\fancyfoot[C]{\small\sffamily\hyperref[nav:inicio]{Índice}}
\fancyfoot[R]{\small\sffamily\bfseries\color{Grafito}Página~\thepage\ de~\pageref*{LastPage}}
\renewcommand{\headrulewidth}{0.35pt}
\renewcommand{\footrulewidth}{0.35pt}
\titleformat{\section}{\sffamily\LARGE\bfseries\color{Grafito}}
  {\textcolor{Enlace}{\thesection}}{0.65em}{}
\titleformat{\subsection}{\sffamily\large\bfseries\color{Grafito}}
  {\thesubsection}{0.65em}{}
\titlespacing*{\section}{0pt}{0pt}{10pt}
\titlespacing*{\subsection}{0pt}{9pt}{2pt}

\newcommand{\Completar}[1]{\textcolor{Gris}{\textit{[#1]}}}
\newcommand{\Guia}[2]{%
  \begin{tcolorbox}[colback=Arena,colframe=Terracota,boxrule=0.35pt,
    leftrule=2pt,arc=0pt,left=7pt,right=7pt,
    top=4pt,bottom=4pt,before skip=5pt,after skip=5pt]
    {\sffamily\small\bfseries\color{Grafito}#1}\par #2
  \end{tcolorbox}}
\newcommand{\EncTabla}[1]{\textcolor{white}{\sffamily\small\bfseries #1}}
\newcommand{\CabeceraTabla}{\rowcolor{Grafito}}
\newcommand{\Capitulo}[4]{%
  \clearpage\section{#1}\label{#2}%
  \noindent{\sffamily\small
    \hyperref[#3]{\(\leftarrow\) Anterior}\enspace\textcolor{Gris}{|}
    \enspace\hyperref[nav:inicio]{Índice}\enspace\textcolor{Gris}{|}
    \enspace\hyperref[sec:catalogo]{Registro de evidencias}%
    \enspace\textcolor{Gris}{|}\enspace
    \hyperref[#4]{Siguiente \(\rightarrow\)}}\par\smallskip}
\newcommand{\Cierre}[2]{%
  \Guia{Decisión y transferencia al siguiente capítulo}{%
    \Completar{Enuncien la decisión, sus límites y el dato o documento
    vigente que entregan para la siguiente etapa.}
    Continúen en \hyperref[#1]{#2}.}}

% B. REGISTRO CENTRAL DE EVIDENCIAS
\newcommand{\ListaEvidencias}{}
\newcommand{\Registrar}[6]{%
  \csdef{evnombre:#1}{#2}\csdef{evcap:#1}{#3}%
  \csdef{evcarpeta:#1}{#4}\csdef{evrev:#1}{#5}\csdef{evfecha:#1}{#6}%
  \gappto\ListaEvidencias{\FilaEvidencia{#1}}}
\NewDocumentCommand{\URL}{m v}{\csgdef{evurl:#1}{#2}}
\newcommand{\NombreEv}[1]{\csuse{evnombre:#1}}
\newcommand{\Ficha}[1]{\hyperlink{ev:#1}{\textsf{\textbf{#1}}}}
\newcommand{\AlCapitulo}[1]{\hyperref[\csuse{evcap:#1}]{§\,\ref*{\csuse{evcap:#1}}}}
\newcommand{\Abrir}[3][URL pendiente]{%
  \ifcsdef{evurl:#2}{\href{\csuse{evurl:#2}}{#3}}
    {\textcolor{Gris}{\textit{#1}}}}
\newcommand{\Acceso}[1]{%
  \par\noindent{\small\Ficha{#1}\enspace\NombreEv{#1}\enspace
  \textcolor{Gris}{|}\enspace\Abrir{#1}{Abrir evidencia}}\par}
\newcommand{\AccesoRaiz}{%
  \ifcsdef{evurl:DRIVE}{\href{\csuse{evurl:DRIVE}}{Carpeta raíz de Drive}}
  {\hyperref[sec:drive]{Carpeta de Drive · configurar enlace}}}

% Incluya una vista legible en Overleaf. El archivo original y su proyecto
% editable se enlazan desde Drive. Es posible compilar sin las vistas.
\newcommand{\Plano}[4]{%
  \par\smallskip
  \begin{minipage}{\linewidth}
  \centering
  \IfFileExists{#1}{%
    \includegraphics[width=\linewidth,height=93mm,keepaspectratio]{#1}}{%
    \begin{tcolorbox}[colback=Papel,colframe=Separador,arc=0pt,
      boxrule=0.35pt,halign=center,top=4pt,bottom=4pt]
      \small\sffamily\color{Gris}Insertar vista legible de #2
    \end{tcolorbox}}%
  \captionof{figure}{#2}\label{#3}%
  {\small\Ficha{#4}\enspace\Abrir{#4}{Abrir plano en Drive}\par}
  \end{minipage}\par
}
\captionsetup{font=small,labelfont={sf,bf},skip=4pt,hypcap=false}

% 01-Firma-y-bases · cap:bases
\Registrar{E0}{Conformación de la firma de ingeniería}{cap:bases}{01-Firma-y-bases}{v1.0}{17/08/2026}
\Registrar{RFP}{Solicitud de propuesta}{cap:bases}{01-Firma-y-bases}{v1.0}{14/08/2026}
\Registrar{BD-PRE-01}{Bases preliminares de diseño}{cap:bases}{01-Firma-y-bases}{v1.0}{25/08/2026}
\Registrar{A2.1}{Matriz de bases de diseño}{cap:bases}{01-Firma-y-bases}{v1.0}{25/08/2026}
\Registrar{NB-01}{Notebook de balance estequiométrico}{cap:bases}{01-Firma-y-bases}{v1.0}{25/08/2026}
\Registrar{NB01-COLAB}{Apertura del notebook NB-01 en Colab}{cap:bases}{01-Firma-y-bases}{v1.0}{25/08/2026}
\Registrar{NB01-INSUMOS}{Datos y auxiliares de NB-01}{cap:bases}{01-Firma-y-bases}{v1.0}{25/08/2026}

% 01-Firma-y-bases · cap:conceptual
\Registrar{AR-01}{Revisión de la propuesta conceptual}{cap:conceptual}{01-Firma-y-bases}{v1.0}{28/08/2026}
\Registrar{EQ-01}{Propuesta preliminar de equipos}{cap:conceptual}{01-Firma-y-bases}{v1.0}{28/08/2026}
\Registrar{E1}{Bases de diseño consolidadas}{cap:conceptual}{01-Firma-y-bases}{v1.0}{04/09/2026}

% 03-AutoCAD-Plant3D · cap:conceptual
\Registrar{DB-V0}{DB/BFD v0 en PDF}{cap:conceptual}{03-AutoCAD-Plant3D}{v0.0}{24/08/2026}
\Registrar{DB-V0-CAD}{Proyecto completo de Plant 3D que respalda DB/BFD v0}{cap:conceptual}{03-AutoCAD-Plant3D}{v0.0}{24/08/2026}

% 02-DWSIM · cap:dwsim
\Registrar{DWSIM-V0}{DWSIM v0, paquete editable completo}{cap:dwsim}{02-DWSIM}{v0.0}{25/08/2026}
\Registrar{DWSIM-V1}{DWSIM v1, paquete editable completo}{cap:dwsim}{02-DWSIM}{v1.0}{11/09/2026}
\Registrar{SIM-01}{Reporte de simulación y verificación v0–v1}{cap:dwsim}{02-DWSIM}{v1.0}{11/09/2026}
\Registrar{BALANCES}{Memorias de materia, energía, reacción y separación}{cap:dwsim}{02-DWSIM}{v1.0}{11/09/2026}
\Registrar{SENS-TEC}{Resultados de sensibilidad técnica}{cap:dwsim}{02-DWSIM}{v1.0}{14/09/2026}

% 03-AutoCAD-Plant3D · cap:dwsim
\Registrar{DB-V1}{DB/BFD v1 en PDF}{cap:dwsim}{03-AutoCAD-Plant3D}{v1.0}{04/09/2026}
\Registrar{DB-V1-CAD}{Proyecto completo de Plant 3D que respalda DB/BFD v1}{cap:dwsim}{03-AutoCAD-Plant3D}{v1.0}{04/09/2026}

% 02-DWSIM · cap:integrado
\Registrar{DWSIM-V2}{DWSIM v2, paquete del caso base integrado}{cap:integrado}{02-DWSIM}{v2.0}{28/09/2026}
\Registrar{VTP-01}{Verificación técnica del proceso}{cap:integrado}{02-DWSIM}{v1.0}{21/09/2026}
\Registrar{TC-01}{Tabla de corrientes de materia}{cap:integrado}{02-DWSIM}{v1.0}{21/09/2026}
\Registrar{TE-01}{Tabla de cargas térmicas, potencias y energía}{cap:integrado}{02-DWSIM}{v1.0}{21/09/2026}
\Registrar{LE-01}{Lista maestra de equipos y variables de diseño}{cap:integrado}{02-DWSIM}{v1.0}{24/09/2026}
\Registrar{HD-EQUIPOS}{Hojas de datos de equipos}{cap:integrado}{02-DWSIM}{v1.0}{24/09/2026}
\Registrar{SERVICIOS}{Balance de servicios e integración energética}{cap:integrado}{02-DWSIM}{v1.0}{05/10/2026}

% 03-AutoCAD-Plant3D · cap:integrado
\Registrar{DFP-V0}{DFP/PFD v0 en PDF}{cap:integrado}{03-AutoCAD-Plant3D}{v0.0}{21/09/2026}
\Registrar{DFP-V0-CAD}{Proyecto completo de Plant 3D que respalda DFP v0}{cap:integrado}{03-AutoCAD-Plant3D}{v0.0}{21/09/2026}
\Registrar{DFP-V1}{DFP/PFD v1 en PDF}{cap:integrado}{03-AutoCAD-Plant3D}{v1.0}{24/09/2026}
\Registrar{DFP-V1-CAD}{Proyecto completo de Plant 3D que respalda DFP v1}{cap:integrado}{03-AutoCAD-Plant3D}{v1.0}{24/09/2026}

% 04-Ambiente-y-SST · cap:ambiente
\Registrar{AC-AMB-01}{Matriz integrada de aspectos ambientales}{cap:ambiente}{04-Ambiente-y-SST}{v1.0}{21/09/2026}
\Registrar{AC-SST-01}{Expediente de seguridad y salud en el trabajo}{cap:ambiente}{04-Ambiente-y-SST}{v1.0}{08/10/2026}
\Registrar{SST-A}{Identificación y valoración del riesgo inicial}{cap:ambiente}{04-Ambiente-y-SST}{v1.0}{08/10/2026}
\Registrar{SST-B}{Requisitos normativos y controles}{cap:ambiente}{04-Ambiente-y-SST}{v1.0}{08/10/2026}
\Registrar{SST-C}{Riesgo residual y decisiones}{cap:ambiente}{04-Ambiente-y-SST}{v1.0}{08/10/2026}

% 05-Control-y-HAZOP · cap:control
\Registrar{FC-01}{Filosofía de control y matriz maestra}{cap:control}{05-Control-y-HAZOP}{v1.0}{08/10/2026}
\Registrar{AC-HAZOP-01}{Estudio HAZOP por nodos y acciones}{cap:control}{05-Control-y-HAZOP}{v1.0}{19/10/2026}
\Registrar{HZ-A}{Formato HZ-A}{cap:control}{05-Control-y-HAZOP}{v1.0}{19/10/2026}
\Registrar{HZ-B}{Formato HZ-B}{cap:control}{05-Control-y-HAZOP}{v1.0}{19/10/2026}
\Registrar{MCE-01}{Matriz causa–efecto}{cap:control}{05-Control-y-HAZOP}{v1.0}{19/10/2026}
\Registrar{CTRL-IPYNB}{Notebook ejecutado de controladores en Python}{cap:control}{05-Control-y-HAZOP}{v1.0}{28/10/2026}
\Registrar{CTRL-COLAB}{Apertura del notebook de controladores en Colab}{cap:control}{05-Control-y-HAZOP}{v1.0}{28/10/2026}
\Registrar{CTRL-INSUMOS}{Datos, scripts auxiliares y versiones de bibliotecas}{cap:control}{05-Control-y-HAZOP}{v1.0}{28/10/2026}
\Registrar{RT-03-CTRL}{Correspondencia de lazos, TAG, escenarios y métricas}{cap:control}{05-Control-y-HAZOP}{v1.0}{28/10/2026}

% 03-AutoCAD-Plant3D · cap:control
\Registrar{DTI-V0}{DTI/P\&ID v0 en PDF}{cap:control}{03-AutoCAD-Plant3D}{v0.0}{08/10/2026}
\Registrar{DTI-V0-CAD}{Proyecto completo de Plant 3D que respalda DTI v0}{cap:control}{03-AutoCAD-Plant3D}{v0.0}{08/10/2026}
\Registrar{DTI-V1}{DTI/P\&ID v1 en PDF}{cap:control}{03-AutoCAD-Plant3D}{v1.0}{19/10/2026}
\Registrar{DTI-V1-CAD}{Proyecto completo de Plant 3D que respalda DTI v1}{cap:control}{03-AutoCAD-Plant3D}{v1.0}{19/10/2026}

% 06-PIB · cap:pib
\Registrar{PIB-01}{Paquete de ingeniería básica y memoria}{cap:pib}{06-PIB}{v1.0}{28/10/2026}
\Registrar{F-PIB-01}{Bases consolidadas}{cap:pib}{06-PIB}{v1.0}{28/10/2026}
\Registrar{F-PIB-02}{Inventario técnico de equipos}{cap:pib}{06-PIB}{v1.0}{28/10/2026}
\Registrar{F-PIB-03}{Balance de recursos, servicios y salidas}{cap:pib}{06-PIB}{v1.0}{28/10/2026}
\Registrar{F-PIB-04}{Pendientes y verificación}{cap:pib}{06-PIB}{v1.0}{28/10/2026}
\Registrar{AG-01}{Arreglo general}{cap:pib}{06-PIB}{v1.0}{28/10/2026}
\Registrar{HAZOP-RES}{Resumen del HAZOP incorporado al PIB}{cap:pib}{06-PIB}{v1.0}{28/10/2026}

% 07-BioSTEAM-y-economia · cap:costos
\Registrar{BIOSTEAM}{Notebook ejecutado de BioSTEAM}{cap:costos}{07-BioSTEAM-y-economia}{v1.0}{05/11/2026}
\Registrar{BIO-COLAB}{Apertura del notebook de BioSTEAM en Colab}{cap:costos}{07-BioSTEAM-y-economia}{v1.0}{05/11/2026}
\Registrar{BIO-INSUMOS}{Datos, módulos, dependencias y versión de bibliotecas}{cap:costos}{07-BioSTEAM-y-economia}{v1.0}{05/11/2026}
\Registrar{RT-03-MAP}{Comparación DWSIM–BioSTEAM}{cap:costos}{07-BioSTEAM-y-economia}{v1.0}{05/11/2026}
\Registrar{RT-03-FE}{Fuentes y bases para el costeo}{cap:costos}{07-BioSTEAM-y-economia}{v1.0}{05/11/2026}
\Registrar{RT-03-CAPEX}{Memoria y tabla de inversión por equipo y total}{cap:costos}{07-BioSTEAM-y-economia}{v1.0}{05/11/2026}
\Registrar{RT-03-OPEX}{Memoria y tabla de costos fijos y variables}{cap:costos}{07-BioSTEAM-y-economia}{v1.0}{05/11/2026}

% 07-BioSTEAM-y-economia · cap:tea
\Registrar{TEA}{Modelo y memoria de evaluación tecnoeconómica}{cap:tea}{07-BioSTEAM-y-economia}{v1.0}{14/11/2026}
\Registrar{FLUJO-CAJA}{Flujo de caja editable e indicadores}{cap:tea}{07-BioSTEAM-y-economia}{v1.0}{14/11/2026}
\Registrar{SENS-ECO}{Escenarios, sensibilidad y resultados reproducibles}{cap:tea}{07-BioSTEAM-y-economia}{v1.0}{14/11/2026}

% 08-Revisiones-tecnicas · cap:tea
\Registrar{PITCH}{Presentación o video de la propuesta final}{cap:tea}{08-Revisiones-tecnicas}{v1.0}{20/11/2026}

% 08-Revisiones-tecnicas · sec:registro
\Registrar{RT-01}{Carpeta y expediente de la revisión técnica 1}{sec:registro}{08-Revisiones-tecnicas}{v1.0}{28/08/2026}
\Registrar{RT01-MATERIAL}{Presentación o evidencia solicitada en RT-01}{sec:registro}{08-Revisiones-tecnicas}{v1.0}{28/08/2026}
\Registrar{RT01-OBS}{Retroalimentación y acciones de RT-01}{sec:registro}{08-Revisiones-tecnicas}{v1.0}{28/08/2026}
\Registrar{RT-02}{Carpeta y expediente de la revisión técnica 2}{sec:registro}{08-Revisiones-tecnicas}{v1.0}{28/09/2026}
\Registrar{RT02-MATERIAL}{Presentación editable y material de RT-02}{sec:registro}{08-Revisiones-tecnicas}{v1.0}{28/09/2026}
\Registrar{RT02-VIDEO}{Video de RT-02}{sec:registro}{08-Revisiones-tecnicas}{v1.0}{28/09/2026}
\Registrar{RT02-OBS}{Retroalimentación y acciones de RT-02}{sec:registro}{08-Revisiones-tecnicas}{v1.0}{28/09/2026}
\Registrar{RT-03}{Carpeta y expediente de la revisión técnica 3}{sec:registro}{08-Revisiones-tecnicas}{v1.0}{14/11/2026}
\Registrar{RT03-MATERIAL}{Presentación editable y material de RT-03}{sec:registro}{08-Revisiones-tecnicas}{v1.0}{14/11/2026}
\Registrar{RT03-VIDEO}{Video de RT-03}{sec:registro}{08-Revisiones-tecnicas}{v1.0}{14/11/2026}
\Registrar{RT-03-RO}{Registro de observaciones, acciones y cierres}{sec:registro}{08-Revisiones-tecnicas}{v1.0}{14/11/2026}

% 00-Portafolio · sec:registro
\Registrar{EF}{Paquete final con PDF, código LaTeX y recursos}{sec:registro}{00-Portafolio}{Rev. 1.0}{23/11/2026}

% C. CALENDARIO Y GANTT
\newcommand{\AgendaPrimera}{%
  \Evento{E0}{Firma de ingeniería}{1}{1}{1}{17/08/2026}{—}{E}
  \Evento{RFP}{Requisitos de la propuesta}{1}{1}{1}{20/08/2026}{E0}{E}
  \Evento{A2.1}{Matriz de bases de diseño}{2}{2}{2}{25/08/2026}{RFP}{E}
  \Evento{NB-01}{Balance estequiométrico}{2}{2}{2}{25/08/2026}{A2.1}{E}
  \Evento{DB-V0}{DB/BFD v0 y proyecto CAD}{2}{2}{2}{24/08/2026}{NB-01}{E}
  \Evento{BD-PRE-01}{Bases preliminares}{2}{2}{2}{25/08/2026}{DB-V0}{E}
  \Evento{AR-01}{Revisión conceptual}{2}{2}{2}{28/08/2026}{BD-PRE-01}{E}
  \Evento{EQ-01}{Propuesta de equipos}{3}{3}{3}{28/08/2026}{AR-01}{E}
  \Evento{DWSIM-V0}{Simulación exploratoria}{3}{3}{3}{01/09/2026}{EQ-01}{E}
  \Evento{DB-V1}{DB/BFD v1 y proyecto CAD}{3}{3}{3}{04/09/2026}{DWSIM-V0}{E}
  \Evento{E1}{Bases consolidadas}{2}{3}{3}{04/09/2026}{DB-V1}{E}
  \Evento{SIM-01}{Verificación de DWSIM v1}{3}{4}{4}{11/09/2026}{DWSIM-V0}{E}
  \Evento{DWSIM-V1}{Caso verificado de DWSIM}{4}{4}{4}{11/09/2026}{SIM-01}{E}
  \Evento{SENS-TEC}{Sensibilidad técnica}{4}{4}{4}{14/09/2026}{DWSIM-V1}{E}
  \Evento{VTP-01}{Verificación técnica}{5}{5}{5}{21/09/2026}{DWSIM-V1}{E}
  \Evento{TC-01}{Corrientes de materia}{5}{5}{5}{21/09/2026}{VTP-01}{E}
  \Evento{TE-01}{Cargas y potencias}{5}{5}{5}{21/09/2026}{VTP-01}{E}
  \Evento{DFP-V0}{DFP v0 y proyecto CAD}{5}{5}{5}{21/09/2026}{TC-01}{E}
  \Evento{LE-01}{Equipos y hojas de datos}{5}{6}{6}{24/09/2026}{DFP-V0}{E}
  \Evento{DFP-V1}{DFP v1 y proyecto CAD}{6}{6}{6}{24/09/2026}{LE-01}{E}
  \Evento{RT-01}{Revisión técnica 1}{6}{6}{6}{28/08/2026}{—}{R}
}
\newcommand{\AgendaSegunda}{%
  \Evento{DWSIM-V2}{Caso base integrado}{6}{7}{7}{28/09/2026}{DFP-V1}{E}
  \Evento{SERVICIOS}{Servicios e integración energética}{7}{8}{8}{05/10/2026}{DWSIM-V2}{E}
  \Evento{AC-AMB-01}{Aspectos ambientales}{8}{9}{9}{21/09/2026}{DFP-V1}{E}
  \Evento{RT-02}{Revisión técnica 2}{9}{9}{9}{28/09/2026}{AC-AMB-01}{R}
  \Evento{AC-SST-01}{Seguridad y salud en el trabajo}{9}{10}{10}{08/10/2026}{AC-AMB-01}{E}
  \Evento{FC-01}{Filosofía de control}{10}{11}{11}{08/10/2026}{AC-SST-01}{E}
  \Evento{DTI-V0}{DTI v0 y proyecto CAD}{11}{11}{11}{08/10/2026}{FC-01}{E}
  \Evento{AC-HAZOP-01}{HAZOP por nodos}{11}{12}{12}{19/10/2026}{DTI-V0}{E}
  \Evento{MCE-01}{Matriz causa–efecto}{12}{12}{12}{19/10/2026}{AC-HAZOP-01}{E}
  \Evento{DTI-V1}{DTI v1 y proyecto CAD}{12}{12}{12}{19/10/2026}{MCE-01}{E}
  \Evento{CTRL-IPYNB}{Notebook de controladores}{11}{13}{13}{28/10/2026}{DTI-V1}{E}
  \Evento{PIB-01}{Paquete y formatos de ingeniería básica}{12}{13}{13}{28/10/2026}{DTI-V1}{E}
  \Evento{BIOSTEAM}{Modelo BioSTEAM verificado}{13}{14}{14}{05/11/2026}{PIB-01}{E}
  \Evento{RT-03-CAPEX}{CAPEX y bases de costeo}{14}{14}{14}{05/11/2026}{BIOSTEAM}{E}
  \Evento{RT-03-OPEX}{OPEX fijo y variable}{14}{14}{14}{05/11/2026}{BIOSTEAM}{E}
  \Evento{RT-03}{Revisión técnica 3}{14}{14}{14}{14/11/2026}{RT-03-OPEX}{R}
  \Evento{TEA}{Indicadores y flujo de caja}{14}{15}{15}{17/11/2026}{RT-03-OPEX}{E}
  \Evento{SENS-ECO}{Sensibilidad y escenarios económicos}{15}{15}{15}{18/11/2026}{TEA}{E}
  \Evento{PITCH}{Propuesta y defensa final}{15}{16}{16}{20/11/2026}{TEA}{E}
  \Evento{EF}{Portafolio final y archivos fuente}{15}{16}{16}{\FechaFinal}{PITCH}{E}
}

\newcommand{\Agenda}{\AgendaPrimera\AgendaSegunda}

\newcommand{\VinculoCalendario}[1]{\hyperlink{cal:#1}{\textsf{#1}}}
\newcommand{\CapturarHito}[8]{%
  \csgdef{hitoev:#1}{#1}\csgdef{semanahito:#1}{#5}}
\begingroup\let\Evento\CapturarHito\Agenda\endgroup
% Asocie los archivos auxiliares al hito que gobierna su entrega.
\newcommand{\AsociarHito}[2]{\csdef{hitoev:#1}{#2}}
\AsociarHito{NB01-COLAB}{NB-01}
\AsociarHito{NB01-INSUMOS}{NB-01}
\AsociarHito{DB-V0-CAD}{DB-V0}
\AsociarHito{DB-V1-CAD}{DB-V1}
\AsociarHito{BALANCES}{SIM-01}
\AsociarHito{HD-EQUIPOS}{LE-01}
\AsociarHito{DFP-V0-CAD}{DFP-V0}
\AsociarHito{DFP-V1-CAD}{DFP-V1}
\AsociarHito{SST-A}{AC-SST-01}
\AsociarHito{SST-B}{AC-SST-01}
\AsociarHito{SST-C}{AC-SST-01}
\AsociarHito{DTI-V0-CAD}{DTI-V0}
\AsociarHito{DTI-V1-CAD}{DTI-V1}
\AsociarHito{HZ-A}{AC-HAZOP-01}
\AsociarHito{HZ-B}{AC-HAZOP-01}
\AsociarHito{CTRL-COLAB}{CTRL-IPYNB}
\AsociarHito{CTRL-INSUMOS}{CTRL-IPYNB}
\AsociarHito{RT-03-CTRL}{CTRL-IPYNB}
\AsociarHito{F-PIB-01}{PIB-01}
\AsociarHito{F-PIB-02}{PIB-01}
\AsociarHito{F-PIB-03}{PIB-01}
\AsociarHito{F-PIB-04}{PIB-01}
\AsociarHito{AG-01}{PIB-01}
\AsociarHito{HAZOP-RES}{PIB-01}
\AsociarHito{BIO-COLAB}{BIOSTEAM}
\AsociarHito{BIO-INSUMOS}{BIOSTEAM}
\AsociarHito{RT-03-MAP}{BIOSTEAM}
\AsociarHito{RT-03-FE}{RT-03-CAPEX}
\AsociarHito{FLUJO-CAJA}{TEA}
\AsociarHito{RT01-MATERIAL}{RT-01}
\AsociarHito{RT01-OBS}{RT-01}
\AsociarHito{RT02-MATERIAL}{RT-02}
\AsociarHito{RT02-VIDEO}{RT-02}
\AsociarHito{RT02-OBS}{RT-02}
\AsociarHito{RT03-MATERIAL}{RT-03}
\AsociarHito{RT03-VIDEO}{RT-03}
\AsociarHito{RT-03-RO}{RT-03}
\newcommand{\SemanaEvidencia}[1]{%
  \ifcsdef{hitoev:#1}{\hyperlink{cal:\csuse{hitoev:#1}}{%
    S\csuse{semanahito:\csuse{hitoev:#1}}}}{}}
\newcommand{\Antecedente}[1]{%
  \ifcsdef{evnombre:#1}{\VinculoCalendario{#1}}{#1}}
\newcommand{\FilaCalendario}[8]{%
  \hypertarget{cal:#1}{\Ficha{#1}} &
  #2\enspace\AlCapitulo{#1} &
  S#5 &
  #6 &
  \Antecedente{#7} &
  \hyperlink{gantt:#1}{Gantt}\enspace\Abrir[Pend.]{#1}{Drive}\\\hline
}
\newcommand{\CeldaGantt}[6]{%
  \ifnum#1=#4\relax
    \cellcolor{Enlace}\hyperlink{ev:#6}{\textcolor{white}{\sffamily\bfseries#5}}%
  \else
    \ifnum#1<#2\relax\else
      \ifnum#1>#3\relax\else\cellcolor{Terracota!32!white}\fi
    \fi
  \fi
  \rule{0pt}{1.45ex}%
}
\newcommand{\FilaGantt}[8]{%
  \hypertarget{gantt:#1}{\VinculoCalendario{#1}}\enspace #2 &
  \Antecedente{#7} &
  \CeldaGantt{1}{#3}{#4}{#5}{#8}{#1}&\CeldaGantt{2}{#3}{#4}{#5}{#8}{#1}&
  \CeldaGantt{3}{#3}{#4}{#5}{#8}{#1}&\CeldaGantt{4}{#3}{#4}{#5}{#8}{#1}&
  \CeldaGantt{5}{#3}{#4}{#5}{#8}{#1}&\CeldaGantt{6}{#3}{#4}{#5}{#8}{#1}&
  \CeldaGantt{7}{#3}{#4}{#5}{#8}{#1}&\CeldaGantt{8}{#3}{#4}{#5}{#8}{#1}&
  \CeldaGantt{9}{#3}{#4}{#5}{#8}{#1}&\CeldaGantt{10}{#3}{#4}{#5}{#8}{#1}&
  \CeldaGantt{11}{#3}{#4}{#5}{#8}{#1}&\CeldaGantt{12}{#3}{#4}{#5}{#8}{#1}&
  \CeldaGantt{13}{#3}{#4}{#5}{#8}{#1}&\CeldaGantt{14}{#3}{#4}{#5}{#8}{#1}&
  \CeldaGantt{15}{#3}{#4}{#5}{#8}{#1}&\CeldaGantt{16}{#3}{#4}{#5}{#8}{#1}\\\hline
}
\newcommand{\FilaEvidencia}[1]{%
  \hypertarget{ev:#1}{\textsf{\textbf{#1}}} &
  \NombreEv{#1}\par{\footnotesize\color{Gris}\csuse{evcarpeta:#1}} &
  \csuse{evrev:#1}\par\csuse{evfecha:#1} &
  \AlCapitulo{#1}\par\SemanaEvidencia{#1} &
  \Abrir{#1}{Abrir archivo}\\\hline
}

\begin{document}
\color{Tinta}
\pagenumbering{arabic}
\hypersetup{pdftitle={Portafolio final de Diseño de Procesos y Productos Químicos},
            pdfauthor={\Firma},pdfsubject={Expediente de ingeniería y evaluación tecnoeconómica}}

% PORTADA
\thispagestyle{empty}
\pdfbookmark[0]{Portada}{portada}
{\sffamily\small\bfseries\color{Grafito}
  \Institucion\par \Facultad\par}
\vspace{5mm}
{\color{Terracota}\rule{\linewidth}{1.5pt}}\par
\vspace{13mm}
\begin{minipage}[c]{0.65\linewidth}
  {\sffamily\small\bfseries\color{Enlace}\Asignatura\par}
  \vspace{4mm}
  {\sffamily\bfseries\fontsize{30}{33}\selectfont\color{Grafito}
    Portafolio final\par}
  \vspace{4mm}
  {\sffamily\Large\bfseries\Proyecto\par}
  \vspace{2mm}
  {\sffamily\large\Firma\par}
  {\sffamily\small\color{Gris}\SiglasFirma\enspace|\enspace
    Conformación \FechaConformacion\par}
\end{minipage}\hfill
\begin{minipage}[c]{0.28\linewidth}
  \centering
  \IfFileExists{\LogoEmpresa}{%
    \includegraphics[width=\linewidth,height=35mm,keepaspectratio]{\LogoEmpresa}
  }{%
    \begin{tcolorbox}[height=35mm,colback=Papel,colframe=Separador,
      boxrule=0.55pt,arc=0pt,valign=center,halign=center]
      {\sffamily\small\color{Gris}Logotipo de la firma}
    \end{tcolorbox}}
\end{minipage}

\vspace{8mm}
{\sffamily\large\bfseries\color{Grafito}Conformación de la firma de ingeniería\par}
\begin{tabularx}{\linewidth}{@{}Y Y Y@{}}
\CabeceraTabla \EncTabla{Integrante} & \EncTabla{Cargo} & \EncTabla{Responsabilidad}\\
\IntegrantesFirma
\bottomrule
\end{tabularx}\par
{\small Acta y datos completos de constitución \Ficha{E0}\enspace
\Abrir{E0}{Abrir acta en Drive}.}

\vfill
\begin{tcolorbox}[colback=Arena,colframe=Terracota,boxrule=0.4pt,
  leftrule=2pt,arc=0pt]
  {\sffamily\bfseries\Docente\par}
  {\sffamily\small Sección \Grupo\enspace|\enspace Periodo \Periodo\par
    Entrega \FechaFinal\enspace|\enspace Revisión \Revision\par}
\end{tcolorbox}
{\sffamily\small\hyperref[nav:inicio]{Abrir índice \(\rightarrow\)}
  \enspace|\enspace\AccesoRaiz}

% PRELIMINARES
\clearpage
\phantomsection\label{nav:inicio}
\pdfbookmark[0]{Inicio e índice}{inicio}
{\sffamily\LARGE\bfseries\color{Grafito}Inicio e índice\par}
{\sffamily\small
  \hyperref[nav:calendario]{Calendario}\enspace|\enspace
  \hyperref[nav:gantt]{Gantt}\enspace|\enspace
  \hyperref[nav:rt]{Revisiones técnicas}\enspace|\enspace
  \hyperref[sec:drive]{Carpeta y estructura de Drive}\enspace|\enspace
  \AccesoRaiz\par}
\Guia{Resumen ejecutivo}{%
  \Completar{Sinteticen el alcance, capacidad, pureza, ruta de proceso,
  decisión de diseño, CAPEX, OPEX, VAN, TIR y precio mínimo de venta.
  Cierren con la viabilidad del caso base y su principal incertidumbre.}
  Enlacen aquí los resultados en
  \hyperref[cap:costos]{CAPEX y OPEX} y \hyperref[cap:tea]{evaluación tecnoeconómica}.}
\noindent\textbf{Control documental}\par
\begin{tabularx}{\linewidth}{@{}p{.20\linewidth}p{.22\linewidth}Y@{}}
\CabeceraTabla \EncTabla{Revisión} & \EncTabla{Fecha} & \EncTabla{Cambio, origen y sección}\\
\Revision & \FechaFinal & \Completar{Identifiquen el cambio y enlacen la revisión o evidencia.}\\
\bottomrule
\end{tabularx}
\renewcommand{\contentsname}{Contenido}
\tableofcontents

% CALENDARIO
\clearpage
\phantomsection\label{nav:calendario}
\addcontentsline{toc}{section}{Calendario de entregas}
{\sffamily\LARGE\bfseries\color{Grafito}Calendario de entregas\par}
{\sffamily\small \hyperref[nav:inicio]{Índice}\enspace|\enspace
  \hyperref[nav:gantt]{Gantt}\enspace|\enspace
  \hyperref[nav:rt]{Revisiones técnicas}\par}
La planeación se expresa en semanas del curso. Completen las fechas y horas
oficiales y ajusten las semanas orientativas de RT-01, RT-02 y RT-03 en
la sección C del código fuente. Las dos vistas se actualizan juntas.
Cada código abre su ficha. El símbolo § abre el capítulo correspondiente.

\begingroup\small\setlength{\tabcolsep}{3.5pt}
\arrayrulecolor{Separador}\setlength{\arrayrulewidth}{0.3pt}
\let\Evento\FilaCalendario
\begin{longtable}{@{}L{.118\linewidth}L{.295\linewidth}L{.059\linewidth}
                      L{.14\linewidth}L{.152\linewidth}L{.14\linewidth}@{}}
\CabeceraTabla \EncTabla{Código} & \EncTabla{Entregable} & \EncTabla{Sem.} &
\EncTabla{Fecha y hora} & \EncTabla{Antecedente} & \EncTabla{Accesos}\\
\endfirsthead
\CabeceraTabla \EncTabla{Código} & \EncTabla{Entregable} & \EncTabla{Sem.} &
\EncTabla{Fecha y hora} & \EncTabla{Antecedente} & \EncTabla{Accesos}\\
\endhead
\Agenda
\end{longtable}
\endgroup

% GANTT
\newcommand{\TablaGantt}[1]{%
  \par\smallskip\begingroup
  \small\setlength{\tabcolsep}{2.7pt}\setlength{\parskip}{0pt}%
  \arrayrulecolor{Separador}\setlength{\arrayrulewidth}{0.3pt}%
  \renewcommand{\arraystretch}{1}%
  \let\Evento\FilaGantt
  {\centering
  \begin{tabular}{@{}L{.310\linewidth}L{.096\linewidth}
    *{16}{>{\centering\arraybackslash}p{.025\linewidth}}@{}}
  \CabeceraTabla \EncTabla{Entregable} & \EncTabla{Antecedente} &
  \EncTabla{1}&\EncTabla{2}&\EncTabla{3}&\EncTabla{4}&
  \EncTabla{5}&\EncTabla{6}&\EncTabla{7}&\EncTabla{8}&
  \EncTabla{9}&\EncTabla{10}&\EncTabla{11}&\EncTabla{12}&
  \EncTabla{13}&\EncTabla{14}&\EncTabla{15}&\EncTabla{16}\\
  #1
  \end{tabular}\par}
  \endgroup
}
\newcommand{\PieGantt}{%
  \vfill\noindent\begin{minipage}{\linewidth}
    \setlength{\parskip}{0pt}%
    {\color{Separador}\rule{\linewidth}{0.4pt}}\par
    {\small\sffamily\color{Grafito}
      \Docente\hfill\hyperref[nav:inicio]{Índice}\hfill
      \textbf{Página~\thepage\ de~\pageref*{LastPage}}}\par
  \end{minipage}\par}
\newcommand{\GuiaGantt}{%
  {\small\sffamily
    \hyperref[nav:calendario]{Calendario}\enspace|\enspace
    \hyperref[nav:rt]{Revisiones técnicas}\par}
  {\small Franja clara: elaboración. \textbf{E}: entrega. \textbf{R}: revisión.
    El código abre el calendario y el marcador abre la ficha de evidencia.\par}}
\begin{landscape}
\thispagestyle{empty}
\enlargethispage{31pt}
\phantomsection\label{nav:gantt}
\addcontentsline{toc}{section}{Gantt de elaboración y entrega}
{\sffamily\LARGE\bfseries\color{Grafito}Gantt de elaboración y entrega\par}
\GuiaGantt
\TablaGantt{\AgendaPrimera}
\PieGantt

\newpage
\thispagestyle{empty}
\enlargethispage{31pt}
{\sffamily\LARGE\bfseries\color{Grafito}Gantt de elaboración y entrega · continuación\par}
\GuiaGantt
\TablaGantt{\AgendaSegunda}
\PieGantt
\end{landscape}


\Capitulo{Firma de ingeniería, alcance y bases preliminares}
  {cap:bases}{nav:calendario}{cap:conceptual}
\subsection{Identidad de la firma y requisitos}
\Completar{Presenten la firma, su organización, sus responsabilidades y
el encargo de la RFP. Definan los límites técnicos, las corrientes de
entrada y salida y las exclusiones acordadas con el cliente.}
\Acceso{E0}
\Acceso{RFP}
\subsection{Matriz de bases y balance preliminar}\label{sec:bases-matriz}
\Completar{Fijen la base de cálculo, capacidad, pureza, horas de operación,
especificaciones, restricciones y supuestos. Incluyan las reacciones
balanceadas, el consumo teórico de reactivos y la producción estimada.
Para electrólisis, documenten agua, electricidad, H\textsubscript{2}
y O\textsubscript{2} sobre la misma base.}\par
\begin{tabularx}{\linewidth}{@{}Y Y Y Y@{}}
\CabeceraTabla \EncTabla{Requisito o variable} & \EncTabla{Valor y unidad} &
  \EncTabla{Tipo y fuente} & \EncTabla{Uso en el diseño}\\
\Completar{Variable} & \Completar{Valor} &
  \Completar{Requisito o supuesto} & \Completar{Decisión asociada}\\
\bottomrule
\end{tabularx}
\Acceso{BD-PRE-01}
\Acceso{A2.1}
\Acceso{NB-01}
\Acceso{NB01-COLAB}
\Acceso{NB01-INSUMOS}
\Guia{Discusión de las bases}{%
  \Completar{Relacionen una exigencia de la RFP con un resultado de NB-01.
  Evalúen el efecto de modificar el supuesto más sensible e indiquen qué
  restricciones deberá respetar el diseño conceptual.}}
\Cierre{cap:conceptual}{Diseño conceptual y equipos}

\Capitulo{Diseño conceptual y selección preliminar de equipos}
  {cap:conceptual}{cap:bases}{cap:dwsim}
\subsection{Traducción de las bases al diagrama de bloques}
Partan de la \hyperref[sec:bases-matriz]{matriz de bases y el balance
preliminar}. \Completar{Justifiquen el orden de los bloques y ubiquen
alimentaciones, producto, subproductos, reciclos y purgas que apliquen.}
\Plano{db-bfd-v0.pdf}{DB/BFD v0}{fig:dbv0}{DB-V0}
\Acceso{DB-V0-CAD}
\subsection{Revisión conceptual y propuesta de equipos}
\Completar{Registren las preguntas y observaciones de AR-01. Vinculen
bloque, TAG, equipo, función, variable por determinar y método para
calcularla en EQ-01. Identifiquen los supuestos que se consolidan en E1.}
\Acceso{AR-01}
\Acceso{EQ-01}
\Acceso{E1}
\Guia{Discusión de las alternativas}{%
  \Completar{Contrasten la configuración preliminar con la alternativa
  que discutieron. Expliquen qué observación de AR-01 cambió la selección
  o el orden de un bloque, qué sustenta el cambio y qué falta verificar.}}
\Cierre{cap:dwsim}{Simulación en DWSIM}

\Capitulo{Simulación y verificación en DWSIM}
  {cap:dwsim}{cap:conceptual}{cap:integrado}
\subsection{Caso exploratorio y cambio de BFD}
Partan del \hyperref[fig:dbv0]{DB/BFD v0} y de \Ficha{EQ-01}.\par
\Completar{Identifiquen el entorno, los componentes, el paquete
termodinámico y las especificaciones de los equipos. Expliquen qué
hallazgos de v0 motivaron cambios en el diagrama de bloques.}
\Acceso{DWSIM-V0}
\Plano{db-bfd-v1.pdf}{DB/BFD v1}{fig:dbv1}{DB-V1}
\Acceso{DB-V1-CAD}
\subsection{SIM-01 y caso verificado}\label{sec:sim-verif}
\Completar{Documenten reacción, separación, reciclo, purga, convergencia
y cierres de materia y energía. Comparen capacidad, pureza y variables
críticas con E1. Muestren rangos y resultados de sensibilidad técnica.}
\begin{tabularx}{\linewidth}{@{}Y Y Y Y@{}}
\CabeceraTabla \EncTabla{Variable y unidad} & \EncTabla{v0} &
\EncTabla{v1 y desviación} & \EncTabla{Criterio y juicio}\\
\Completar{Variable} & \Completar{Valor} & \Completar{Valor y cambio}
  & \Completar{Base o tolerancia}\\
\bottomrule
\end{tabularx}
\Acceso{DWSIM-V1}
\Acceso{SIM-01}
\Acceso{BALANCES}
\Acceso{SENS-TEC}
\Guia{Discusión de los resultados de simulación}{%
  \Completar{Interpreten un cambio entre v0 y v1. Justifiquen si surge
  de la configuración, del paquete termodinámico o de los supuestos.
  Expliquen qué resultados permiten elegir las condiciones para el caso integrado.}}
\Cierre{cap:integrado}{Proceso integrado, tablas y servicios}

\Capitulo{Proceso integrado, DFP, equipos y servicios}
  {cap:integrado}{cap:dwsim}{cap:ambiente}
\subsection{Caso de referencia y tablas de ingeniería}\label{sec:tablas-proceso}
Del \hyperref[sec:sim-verif]{caso verificado} obtengan un caso integrado.
\Completar{Identifiquen la versión y la condición de operación usada
para emitir el DFP. Registren caudal, composición, temperatura, presión
y fase en TC-01. En TE-01, registren cargas y potencias con unidades,
signos y equipo asociado. Contrasten los resultados con VTP-01.}
\Acceso{DWSIM-V2}
\Acceso{VTP-01}
\Acceso{TC-01}
\Acceso{TE-01}
\subsection{Emisión del DFP y consolidación de equipos}
\Completar{Identifiquen los cambios entre DFP v0 y v1. Relacionen los
TAG del plano con LE-01 y las hojas de datos. Indiquen cantidad,
función, tipo, parámetro de tamaño y origen de cada dato.}
\Acceso{DFP-V0}
\Acceso{DFP-V0-CAD}
\Plano{dfp-v1.pdf}{DFP/PFD v1}{fig:dfpv1}{DFP-V1}
\Acceso{DFP-V1-CAD}
\Acceso{LE-01}
\Acceso{HD-EQUIPOS}
\subsection{Servicios e integración energética}\label{sec:servicios}
\Completar{Cuantifiquen demandas y recuperación de calor, cuando
corresponda. Relacionen TE-01 con el consumo externo y con las condiciones
de referencia que se usarán en F-PIB-03.}
\Acceso{SERVICIOS}
\Guia{Discusión de diseño y consumo}{%
  \Completar{Expliquen cómo un resultado del caso integrado cambió
  una especificación de equipo o un servicio. Comparen las alternativas
  evaluadas y justifiquen el caso transferido a ambiente y control.}}
\Cierre{cap:ambiente}{Ambiente y seguridad laboral}

\Capitulo{Aspectos ambientales y seguridad y salud en el trabajo}
  {cap:ambiente}{cap:integrado}{cap:control}
\subsection{Evaluación ambiental integrada}\label{sec:amb-matriz}
Usen el \hyperref[fig:dfpv1]{DFP vigente} y el
\hyperref[sec:servicios]{balance de servicios}.
\Completar{Sinteticen la matriz única AC-AMB-01 con aspecto, impacto,
significancia, medida de control, normatividad y requisito concreto.
Identifiquen la etapa, corriente o TAG de los aspectos prioritarios.}
\Acceso{AC-AMB-01}
{\small \hyperref[nav:rt02]{Consultar RT-02 y sus observaciones}.}
\subsection{Peligros, requisitos y controles}
\Completar{Vinculen los formatos A, B y C por identificador de escenario.
Documenten el peligro, riesgo inicial, requisito aplicable, controles,
riesgo residual y decisión. Enlacen el DFP o el equipo afectado.}
\Acceso{AC-SST-01}
\Acceso{SST-A}
\Acceso{SST-B}
\Acceso{SST-C}
\begin{tabularx}{\linewidth}{@{}Y Y Y@{}}
\CabeceraTabla \EncTabla{Aspecto o escenario} & \EncTabla{Evidencia y criterio} &
\EncTabla{Control y capítulo afectado}\\
\Completar{ID, TAG y corriente} & \Completar{Fila, requisito y medida} &
  \Completar{Decisión y enlace interno}\\
\bottomrule
\end{tabularx}
\Guia{Discusión ambiental y de SST}{%
  \Completar{Expliquen cuál es el aspecto más significativo y cuál es
  el escenario prioritario. Analicen cómo el control propuesto modifica
  el riesgo o el impacto y qué hallazgo se incorpora al diseño o a FC-01.}}
\Cierre{cap:control}{Control, HAZOP y DTI}

\Capitulo{Filosofía de control, HAZOP y DTI}
  {cap:control}{cap:ambiente}{cap:pib}
\subsection{Filosofía de control y representación preliminar}
Recuperen los \hyperref[cap:ambiente]{hallazgos de ambiente y SST}
y el \hyperref[fig:dfpv1]{DFP vigente}.
\Completar{Definan por equipo la variable controlada, consigna o
intervalo, elemento final, perturbación, respuesta, alarmas, interbloqueos
y TAG con su fuente. Expliquen las interacciones relevantes entre lazos.}
\Acceso{FC-01}
\Acceso{DTI-V0}
\Acceso{DTI-V0-CAD}
\subsection{HAZOP, matriz causa--efecto y revisión del DTI}\label{sec:hazop}
\Completar{Expliquen cómo se eligieron los nodos, las desviaciones,
causas, consecuencias, salvaguardas y acciones. Vinculen escenarios,
TAG, condiciones de activación y cambios de plano.}
\Acceso{AC-HAZOP-01}
\Acceso{HZ-A}
\Acceso{HZ-B}
\Acceso{MCE-01}
\Plano{dti-v1.pdf}{DTI/P\&ID v1}{fig:dtiv1}{DTI-V1}
\Acceso{DTI-V1-CAD}
\subsection{Pruebas de los controladores}\label{sec:ctrl}
\Completar{Indiquen los lazos programados, ecuaciones, parámetros de
sintonía, pruebas, gráficas, métricas y conclusiones. Identifiquen la
condición del modelo y el TAG que corresponde a cada lazo.}
\Acceso{CTRL-IPYNB}
\Acceso{CTRL-COLAB}
\Acceso{CTRL-INSUMOS}
\Acceso{RT-03-CTRL}
\Guia{Discusión de la respuesta operativa y las salvaguardas}{%
  \Completar{Sigan un escenario desde el riesgo hasta su acción en la
  matriz causa--efecto y su TAG en el DTI. Expliquen qué prueba y métrica
  del notebook sustentan la estrategia de control y bajo qué límites.}}
\Cierre{cap:pib}{Paquete de ingeniería básica}

\Capitulo{Paquete de ingeniería básica}
  {cap:pib}{cap:control}{cap:costos}
\subsection{Bases, equipos, recursos y verificación}\label{sec:pib-formatos}
Consoliden \hyperref[sec:tablas-proceso]{las tablas de ingeniería},
\hyperref[sec:servicios]{los servicios}, \hyperref[sec:hazop]{el HAZOP}
y \hyperref[fig:dtiv1]{el DTI vigente} en el PIB-01.
\Acceso{PIB-01}
\begin{tabularx}{\linewidth}{@{}p{.19\linewidth}Y Y@{}}
\CabeceraTabla \EncTabla{Formato} & \EncTabla{Contenido que se discute} &
  \EncTabla{Uso posterior}\\
\Ficha{F-PIB-01} & Capacidad, producción, horas y límites de diseño. &
  Base del modelo económico y de la producción vendible.\\
\Ficha{F-PIB-02} & Inventario por TAG, cantidad y variable de tamaño. &
  Cálculo del CAPEX por equipo.\\
\Ficha{F-PIB-03} & Recursos, servicios y salidas por hora y por año. &
  OPEX variable e ingresos.\\
\Ficha{F-PIB-04} & Pendientes, responsables y evidencia de cierre. &
  Supuestos por verificar y sensibilidad.\\
\bottomrule
\end{tabularx}
\Acceso{F-PIB-01}
\Acceso{F-PIB-02}
\Acceso{F-PIB-03}
\Acceso{F-PIB-04}
\subsection{Arreglo general y soportes}
\Completar{Presenten las decisiones de implantación y enlacen la
memoria, el caso base de DWSIM, los planos y las acciones prioritarias.}
\Acceso{AG-01}
\Acceso{HAZOP-RES}
{\small \hyperref[fig:dtiv1]{DTI v1}\enspace|\enspace
\Ficha{DWSIM-V2}\enspace|\enspace\Ficha{TC-01}\enspace|\enspace\Ficha{TE-01}.}
\Guia{Discusión de la consolidación}{%
  \Completar{Sigan el origen de un parámetro de tamaño y de un consumo.
  Identifiquen el dato pendiente de mayor efecto económico y justifiquen
  cómo lo representarán en el análisis siguiente.}}
\Cierre{cap:costos}{BioSTEAM, CAPEX y OPEX}

\Capitulo{BioSTEAM y estimación de CAPEX y OPEX}
  {cap:costos}{cap:pib}{cap:tea}
\subsection{Modelo reproducible y comparación con DWSIM}
Partan del \hyperref[sec:pib-formatos]{PIB-01} y del caso
\Ficha{DWSIM-V2}. \Completar{Describan el entorno, insumos y
alcance de BioSTEAM. Comparen valores de referencia y resultados
con unidades, desviación y justificación. Documenten aproximaciones,
equipos omitidos del modelo y correcciones por fuera de la simulación.}
\Acceso{BIOSTEAM}
\Acceso{BIO-COLAB}
\Acceso{BIO-INSUMOS}
\Acceso{RT-03-MAP}
{\small \hyperref[nav:rt03]{Consultar RT-03 y sus observaciones}.}
\subsection{Inversión}\label{sec:capex}
\Completar{Definan moneda, año monetario, capacidad, horas y límites.
Por TAG indiquen tamaño, cantidad, método y fuente, compra e instalación.
Si usaron Turton/CAPCOST, especifiquen correlación, coeficientes, unidades
y validez. Cuando aplique, muestren el escalamiento y la actualización
por CEPCI. Reconcilien inversión fija, capital de trabajo e inversión total.
Verifiquen la contabilización única de instalación e instrumentación.}
\Acceso{RT-03-FE}
\Acceso{RT-03-CAPEX}
\subsection{Costos operativos}\label{sec:opex}
\Completar{Vinculen F-PIB-03 con consumos, precios, bases y fuentes.
Reporten OPEX fijo, variable, total y costo unitario. Para e-metanol,
especifiquen el costeo de CO\textsubscript{2} y H\textsubscript{2}.
Cuando la electrólisis forme parte de la planta, relacionen el costeo
del hidrógeno con el agua, la electricidad y los otros insumos considerados.}
\Acceso{RT-03-OPEX}
\begin{tabularx}{\linewidth}{@{}Y Y Y Y@{}}
\CabeceraTabla \EncTabla{Partida} & \EncTabla{Base y unidad} &
\EncTabla{Costo y año} & \EncTabla{Origen y cálculo}\\
\Completar{TAG o consumo} & \Completar{Tamaño o flujo} &
  \Completar{Valor} & \Completar{Fila de PIB y memoria}\\
\bottomrule
\end{tabularx}
\Guia{Discusión del costeo y sus límites}{%
  \Completar{Justifiquen cómo se trató la diferencia más relevante
  entre los modelos. Identifiquen el equipo que domina el CAPEX y el
  insumo que domina el OPEX. Expliquen qué supuesto merece sensibilidad.}}
\Cierre{cap:tea}{Evaluación tecnoeconómica}

\Capitulo{Evaluación tecnoeconómica y propuesta final}
  {cap:tea}{cap:costos}{sec:registro}
\subsection{Flujo de caja, supuestos y resultados}\label{sec:flujo}
Integren la \hyperref[sec:capex]{inversión} y los
\hyperref[sec:opex]{costos operativos} del capítulo anterior.
\Completar{Declaren vida del proyecto, tasa de descuento, calendario de
inversión y arranque, producción vendible, precio de venta, depreciación,
impuestos y recuperación de capital de trabajo, según apliquen.
Muestren cómo los datos técnicos se convierten en partidas anuales.}
\Acceso{TEA}
\Acceso{FLUJO-CAJA}
\begin{tabularx}{\linewidth}{@{}p{.33\linewidth}Y Y@{}}
\CabeceraTabla \EncTabla{Indicador} & \EncTabla{Resultado y unidad} &
  \EncTabla{Criterio e interpretación}\\
VAN o VPN & \Completar{Resultado} & \Completar{Tasa y horizonte}\\
TIR & \Completar{Resultado} & \Completar{Tasa de referencia}\\
Periodo de recuperación & \Completar{Años} & \Completar{Simple o descontado}\\
Precio mínimo de venta & \Completar{Moneda por unidad} &
  \Completar{Precio con VAN igual a cero}\\
\bottomrule
\end{tabularx}
\subsection{Sensibilidad y propuesta técnico-económica}
\Completar{Definan caso base, rangos y escenarios. Incorporen resultados
y gráficas. Comparen el precio mínimo de venta con el precio del escenario
y distingan su significado del costo unitario de producción.}
\Acceso{SENS-ECO}
\Acceso{PITCH}
\Guia{Discusión y recomendación de la firma}{%
  \Completar{Concluyan si se cumplen las especificaciones y si el caso
  resulta viable. Describan el escenario que cambiaría la decisión,
  el supuesto más incierto y la verificación necesaria para avanzar.}}
\Cierre{sec:registro}{Revisiones, evidencias y anexos}

\Capitulo{Revisiones técnicas, evidencias y anexos}
  {sec:registro}{cap:tea}{sec:documentacion}
\subsection{Revisiones técnicas}\label{nav:rt}
\begin{itemize}
  \item \hyperref[nav:rt01]{RT-01 · Primer corte técnico}
  \item \hyperref[nav:rt02]{RT-02 · Avances hasta AC-AMB-01}
  \item \hyperref[nav:rt03]{RT-03 · Expediente acumulado, BioSTEAM y costos}
\end{itemize}

\subsubsection*{RT-01 · Primer corte técnico}
\phantomsection\label{nav:rt01}
\Completar{Registren el alcance asignado, versión presentada y fecha.
Vinculen cada observación con el capítulo donde se atendió.}
\Acceso{RT-01}
\Acceso{RT01-MATERIAL}
\Acceso{RT01-OBS}
{\small \hyperlink{cal:RT-01}{Ver fecha en el calendario}.}

\subsubsection*{RT-02 · Avances hasta AC-AMB-01}
\phantomsection\label{nav:rt02}
\Completar{Indiquen qué observaciones modificaron las tablas, el DFP,
la matriz ambiental o la justificación de decisiones.}
\Acceso{RT-02}
\Acceso{RT02-MATERIAL}
\Acceso{RT02-VIDEO}
\Acceso{RT02-OBS}
{\small \hyperlink{cal:RT-02}{Ver fecha}\enspace|\enspace
\hyperref[sec:amb-matriz]{Discusión de AC-AMB-01}.}

\subsubsection*{RT-03 · Expediente acumulado, BioSTEAM y CAPEX/OPEX}
\phantomsection\label{nav:rt03}
\Completar{Señalen las observaciones sobre el expediente, los controladores,
la comparación de simuladores y las tablas de costos. Enlacen las evidencias
de cierre y las limitaciones que se trasladaron a la evaluación económica.}
\Acceso{RT-03}
\Acceso{RT03-MATERIAL}
\Acceso{RT03-VIDEO}
\Acceso{RT-03-RO}
{\small \hyperlink{cal:RT-03}{Ver fecha}\enspace|\enspace
\hyperref[sec:ctrl]{Controladores}\enspace|\enspace
\hyperref[cap:costos]{BioSTEAM y costos}.}

\subsubsection*{Registro de cierre}
\begin{tabularx}{\linewidth}{@{}p{.105\linewidth}Y Y Y@{}}
\CabeceraTabla \EncTabla{Corte} & \EncTabla{Observación y responsable} &
\EncTabla{Acción y estado} & \EncTabla{Evidencia y apartado}\\
\Completar{RT} & \Completar{Hallazgo} & \Completar{Cambio y fecha} &
  \Completar{Ficha y referencia interna}\\
\bottomrule
\end{tabularx}
\Guia{Discusión de la evolución del proyecto}{%
  \Completar{Identifiquen una decisión técnica que cambió con la
  retroalimentación, comparen las versiones anterior y posterior y
  expliquen por qué el cambio mejoró la propuesta o qué queda pendiente.}}

\subsection{Carpeta de Drive y condiciones de verificación}\label{sec:drive}
{\sffamily\bfseries\AccesoRaiz}\par
\Completar{Identifiquen la carpeta por equipo, proyecto y periodo. En el
registro indiquen las revisiones y fechas reales. Comprueben que quien
evaluará el portafolio puede abrir las rutas con su cuenta autorizada.}

\begin{tabularx}{\linewidth}{@{}p{.36\linewidth}Y@{}}
\CabeceraTabla \EncTabla{Subcarpeta} & \EncTabla{Contenido}\\
\texttt{00-Portafolio} & PDF final, código LaTeX, figuras y recursos.\\
\texttt{01-Firma-y-bases} & RFP, bases, notebooks preliminares y revisiones de concepto.\\
\texttt{02-DWSIM} & Subcarpetas \texttt{v0}, \texttt{v1} y \texttt{v2} con cada caso,
  auxiliares, exportaciones y un inventario.\\
\texttt{03-AutoCAD-Plant3D} & Subcarpetas \texttt{DB-BFD}, \texttt{DFP} y
  \texttt{DTI}, con PDF y proyecto completo para cada revisión emitida.\\
\texttt{04-Ambiente-y-SST} & Matriz ambiental y formatos de SST.\\
\texttt{05-Control-y-HAZOP} & FC-01, matrices, notebook y auxiliares.\\
\texttt{06-PIB} & Cuatro formatos, memoria, AG-01 y soportes.\\
\texttt{07-BioSTEAM-y-economia} & Notebook, datos, memorias, costos y flujo de caja.\\
\texttt{08-Revisiones-tecnicas} & Subcarpetas para RT-01, RT-02, RT-03 y pitch.\\
\bottomrule
\end{tabularx}

\begin{itemize}
  \item \textbf{DWSIM.} Guarden, por versión, el archivo de simulación
  y todos los archivos que la acompañan. Identifiquen la versión del software,
  el punto de entrada del caso, dependencias y archivos de respaldo o exportación.
  \item \textbf{Plant 3D.} Conserven un proyecto completo por revisión emitida,
  con sus carpetas y archivos de configuración, dibujos y referencias utilizadas.
  Si un paquete contiene varios planos, vinculen cada plano a ese mismo
  paquete y preserven su revisión. Comprueben que abre desde una copia del paquete.
  \item \textbf{Colab.} Conserven el \texttt{.ipynb} ejecutado, insumos,
  auxiliares y versiones. Enlacen esa copia en Drive y su apertura en Colab.
  \item \textbf{Complementarios.} Registren toda memoria, fuente, hoja de datos
  o archivo usado en la discusión. Se pueden agregar filas en la sección B
  del código y enlaces a ellas en el capítulo correspondiente.
\end{itemize}
\Acceso{EF}

\subsection{Registro maestro de archivos}\label{sec:catalogo}
El código de cada evidencia es el mismo que se usa en sus referencias.
La columna Cap./S enlaza al lugar donde se interpreta el archivo y a
la entrega que le corresponde en el calendario.
Una raya indica que falta completar la revisión o fecha.

\begingroup\small\setlength{\tabcolsep}{4pt}
\arrayrulecolor{Separador}\setlength{\arrayrulewidth}{0.3pt}
\begin{longtable}{@{}L{.174\linewidth}L{.383\linewidth}L{.136\linewidth}
                     L{.085\linewidth}L{.142\linewidth}@{}}
\CabeceraTabla \EncTabla{Código} & \EncTabla{Evidencia y carpeta} &
\EncTabla{Rev. y fecha} & \EncTabla{Cap./S} & \EncTabla{Acceso}\\
\endfirsthead
\CabeceraTabla \EncTabla{Código} & \EncTabla{Evidencia y carpeta} &
\EncTabla{Rev. y fecha} & \EncTabla{Cap./S} & \EncTabla{Acceso}\\
\endhead
\ListaEvidencias
\end{longtable}
\endgroup

\subsection{Referencias y cierre del expediente}\label{sec:referencias}
\Completar{Inserten las referencias citadas. Incluyan fuentes de
propiedades, correlaciones, precios, índices, normatividad y versiones
del software, con fecha de consulta cuando corresponda.}
\Guia{Verificación cruzada antes de entregar}{%
  Desde una \Ficha{LE-01}, sigan el TAG al \hyperref[fig:dfpv1]{DFP},
  a \Ficha{F-PIB-02}, al \hyperref[sec:capex]{CAPEX} y al
  \hyperref[sec:flujo]{flujo de caja}. Desde \Ficha{TE-01}, sigan un
  consumo a \Ficha{F-PIB-03}, al \hyperref[sec:opex]{OPEX} y al mismo
  flujo. Prueben los enlaces de todas las revisiones y las versiones emitidas.
  Contrasten la lista de verificación del curso con este expediente.}

\end{document}
